\section*{Abstract}

% The opening framing follows the paper's abstract (paper/draft/sections/abstract.tex); the results are
% this thesis's. The abstract is citation-free by design; the works it refers to are cited in the body.
Large language models (LLMs) often lie, that is, they generate statements they know or believe to
be false. Detecting these lies is a critical alignment and safety challenge. Prior work demonstrates a
trade-off: computationally expensive black-box approaches achieve higher accuracy, whereas cheaper
white-box methods such as linear probes frequently fail to generalise across datasets. We study
label-free white-box detectors for Llama-3.3-70B-Instruct on the Liars' Bench benchmark. Following the
persona coordinate framework, we build zero-shot directions from honest and deceptive role-play;
extended with roles that deny knowledge, one such direction transfers across the benchmark's datasets.
On lies about knowledge, however, it reads the wording
of a denial rather than the lie: where a lie and an honest answer are the same sentence, no persona
direction beats random directions. We therefore build a second direction from the definition of a lie
itself. The model reads questions it knows and questions nobody can know, each followed by the same
forced denial, and the difference reads what the model knows while the words stay fixed. Without any
labelled lies, this knowledge direction separates lies from identical honest denials (AUROC
\V{th-abs-k-ident}) and detects the lies of a model fine-tuned to hide what it knows
(\V{th-abs-k-gs}, against \V{th-lb-mean-auroc-gs} for the benchmark's supervised probe). Combining the
two gives our selected detector, a single direction built without labelled lies. It reaches a mean
AUROC of \V{th-abs-sel-off} across the benchmark's datasets, against \V{th-abs-base-off} for probes
trained on one dataset and tested on the others, stays ahead of the published persona contrast and of
those probes at every layer we read, and separates lies from identical honest denials. When we add
the persona direction to the model's activations while it answers, the model hides an unfavourable
fact more often; random directions have no such effect. Repeated on the benchmark's three other models,
the persona contrast again transfers better than such probes on Mistral-Small-3.1-24B and Gemma-3-27B but
fails on Qwen2.5-72B, which rarely stays in role; the combined detector beats a probe trained at the same
depth on all four models, but the knowledge direction does not detect the one other fine-tuned liar.

\vspace{1cm}
\textbf{Keywords:} large language models; lie detection; deception; linear probes; persona vectors;
knowledge; model organisms; interpretability; AI safety; out-of-distribution generalisation
\pagebreak
